\documentclass[11pt, a4paper]{article}

% PACKAGES
\usepackage[a4paper, top=2.5cm, bottom=2.5cm, left=2cm, right=2cm]{geometry}
\usepackage{fontspec} % For XeLaTeX
\usepackage[english, bidi=basic, provide=*]{babel}
\usepackage{amsmath}
\usepackage{amssymb}
\usepackage{booktabs}
\usepackage{xcolor}
\usepackage{framed}
\usepackage{fancyvrb}
\usepackage{hyperref}

% PANDOC/CUSTOM DEFINITIONS
\providecommand{\tightlist}{\setlength{\itemsep}{0pt}\setlength{\parskip}{0pt}}

\definecolor{shadecolor}{RGB}{248,248,248}
\newenvironment{Shaded}{\begin{snugshade}}{\end{snugshade}}
\newcommand{\AlertTok}[1]{\textcolor[rgb]{0.94,0.16,0.16}{#1}}
\newcommand{\AttributeTok}[1]{\textcolor[rgb]{0.77,0.63,0.00}{#1}}
\newcommand{\BuiltInTok}[1]{#1}
\newcommand{\CharTok}[1]{\textcolor[rgb]{0.31,0.60,0.02}{#1}}
\newcommand{\CommentTok}[1]{\textcolor[rgb]{0.56,0.35,0.01}{\textit{#1}}}
\newcommand{\DataTypeTok}[1]{\textcolor[rgb]{0.13,0.29,0.53}{#1}}
\newcommand{\DecValTok}[1]{\textcolor[rgb]{0.00,0.00,0.81}{#1}}
\newcommand{\ErrorTok}[1]{\textcolor[rgb]{0.64,0.00,0.00}{\textbf{#1}}}
\newcommand{\ExtensionTok}[1]{#1}
\newcommand{\FloatTok}[1]{\textcolor[rgb]{0.00,0.00,0.81}{#1}}
\newcommand{\FunctionTok}[1]{\textcolor[rgb]{0.00,0.00,0.00}{#1}}
\newcommand{\ImportTok}[1]{#1}
\newcommand{\InformationTok}[1]{\textcolor[rgb]{0.56,0.35,0.01}{\textit{#1}}}
\newcommand{\KeywordTok}[1]{\textcolor[rgb]{0.13,0.29,0.53}{\textbf{#1}}}
\newcommand{\NormalTok}[1]{#1}
\newcommand{\OperatorTok}[1]{\textcolor[rgb]{0.81,0.36,0.00}{\textbf{#1}}}
\newcommand{\OtherTok}[1]{\textcolor[rgb]{0.56,0.35,0.01}{#1}}
\newcommand{\PreprocessorTok}[1]{\textcolor[rgb]{0.56,0.35,0.01}{\textit{#1}}}
\newcommand{\RegionMarkerTok}[1]{#1}
\newcommand{\SpecialCharTok}[1]{\textcolor[rgb]{0.00,0.00,0.00}{#1}}
\newcommand{\SpecialStringTok}[1]{\textcolor[rgb]{0.31,0.60,0.02}{#1}}
\newcommand{\StringTok}[1]{\textcolor[rgb]{0.31,0.60,0.02}{#1}}
\newcommand{\VariableTok}[1]{\textcolor[rgb]{0.00,0.00,0.00}{#1}}
\newcommand{\VerbatimStringTok}[1]{\textcolor[rgb]{0.31,0.60,0.02}{#1}}
\newcommand{\WarningTok}[1]{\textcolor[rgb]{0.56,0.35,0.01}{\textit{#1}}}
\DefineVerbatimEnvironment{Highlighting}{Verbatim}{commandchars=\\\{\}}

\begin{document}

\section{Intention is all you need}\label{intention-is-all-you-need}

\textbf{Author:} Alonso Velazquez\\
\textbf{Affiliation:} Sinaloa Autonomous University\\
\textbf{Date:} March 29, 2026\\
\textbf{arXiv Category:} cs.SE (Software Engineering), cs.PL
(Programming Languages)

\begin{center}\rule{0.5\linewidth}{0.5pt}\end{center}

\subsection{Abstract}\label{abstract}

We introduce \textbf{Semantic Metaprogramming}, a novel paradigm where
natural language intent is treated as a first-class, computable data
type. To demonstrate this, we present \textbf{Glupe}, a proof-of-concept
meta-compiler that enables developers to mix concrete code with
high-level pseudocode or natural language intent within the same source
file. Using a simple delimiter syntax (\texttt{\$\$\{...\}\$\$}),
programmers can write hybrid programs where architectural ``load-bearing
walls'' are specified in traditional code, while algorithmic
implementation details are expressed via human intent.

Glupe transpiles these semantic containers to executable code using
Large Language Models (LLMs), but unlike traditional AI assistants, it
enforces strict architectural sandboxing. We introduce concepts such as
\textbf{Prompt Algebra} and \textbf{Semantic Subtraction}
(\(Spec — Impl\)) to mathematically verify that AI-generated code
conforms to requested constraints without omission. Key features include
targeted dependency injection via a Tree-sitter global symbol graph,
AST-driven structural verification, and iterative semantic equivalence
loops.

\textbf{Keywords:} LLM-assisted programming, hybrid programming
paradigm, semantic metaprogramming, code generation, AST-driven
verification, prompt algebra

\subsection{1. Introduction}\label{introduction}

Large Language Models (LLMs) have shown remarkable ability to generate
code from natural language descriptions. However, current LLM-assisted
programming tools operate primarily through two paradigms:

\begin{enumerate}
\def\labelenumi{\arabic{enumi}.}
\tightlist
\item
  \textbf{Autonomous Agents / Chat interfaces} (Devin, ChatGPT): Full
  program generation that often operates outside the deterministic
  constraints of existing architectures.
\item
  \textbf{Autocomplete} (GitHub Copilot): Inline suggestions that act as
  advanced predictive text but lack formal compilation guarantees.
\end{enumerate}

Both approaches treat natural language as \textbf{external} to the
programming process. We propose a third paradigm: \textbf{Semantic
Metaprogramming}, which treats natural language intent as first-class,
compilable syntax within source files.

\subsubsection{1.1 Motivating Example}\label{motivating-example}

Consider a developer who knows the high-level algorithm but not the
language-specific APIs:

\begin{Shaded}
\begin{Highlighting}[]
\PreprocessorTok{\#include }\ImportTok{\textless{}vector\textgreater{}}
\PreprocessorTok{\#include }\ImportTok{\textless{}iostream\textgreater{}}

\DataTypeTok{int}\NormalTok{ main}\OperatorTok{()} \OperatorTok{\{}
    \ErrorTok{$$}\OperatorTok{\{}
\NormalTok{        create a vector with the first }\DecValTok{10}\NormalTok{ prime numbers}
\NormalTok{        print each one with its index}
\NormalTok{        calculate }\KeywordTok{and}\NormalTok{ print the sum}
    \OperatorTok{\}}\ErrorTok{$$}
    \ControlFlowTok{return} \DecValTok{0}\OperatorTok{;}
\OperatorTok{\}}
\end{Highlighting}
\end{Shaded}

The \texttt{\$\$\{...\}\$\$} blocks are not comments, they are
\textbf{semantic placeholders} compiled to concrete code by an LLM while
preserving the surrounding C++ structure (includes, main signature,
return statement).

This enables: - \textbf{Gradual specification}: Start with intent,
refine to code incrementally - \textbf{Language learning}: Express what
you want, learn syntax from generated output\\
- \textbf{Rapid prototyping}: Focus on architecture, delegate
boilerplate - \textbf{Mixed expertise}: Domain experts collaborate with
language experts in same file

\subsubsection{1.2 Contributions}\label{contributions}

This essay describes the design and implementation of Glupe, a working
meta-compiler supporting this paradigm. Our contributions are:

\begin{enumerate}
\def\labelenumi{\arabic{enumi}.}
\tightlist
\item
  \textbf{Lightweight syntax convention} for embedding semantic blocks
  in source code
\item
  \textbf{Multi-file orchestration} through \texttt{EXPORT} directives
  for project generation
\item
  \textbf{Iterative refinement with error feedback} to improve LLM
  reliability
\item
  \textbf{AST-Driven structural verification and Semantic equivalence
  loops} to catch LLM omissions and anti-patterns
\item
  \textbf{Targeted dependency injection} using a Tree-Sitter based
  global symbol graph
\item
  \textbf{Integrated developer tooling} for code repair, documentation,
  and analysis
\item
  \textbf{Multi-provider LLM backend} supporting local and cloud models
\end{enumerate}

\begin{center}\rule{0.5\linewidth}{0.5pt}\end{center}

\subsection{2. System Design}\label{system-design}

\subsubsection{2.1 Semantic Block Syntax}\label{semantic-block-syntax}

A \textbf{semantic block} is delimited by \texttt{\$\$\{} and
\texttt{\}\$\$} containing natural language or a custom DSL:

\begin{verbatim}
semantic_block ::= "$${" natural_language_text "}$$"
\end{verbatim}

Design rationale: - \textbf{Distinctive markers} (\texttt{\$\$\{},
\texttt{\}\$\$}) avoid conflicts with existing language syntax -
\textbf{Natural language content}: No special formatting required -
\textbf{Position-independent}: Can appear anywhere, code can appear in
target language - \textbf{Identity}: Containers can be named like this:

\begin{verbatim}
$$ my_container { intent goes here }$$
\end{verbatim}

\textbf{Benefits of Semantic Containers}: it effectively sandboxes AI
generation in constrained spaces, but inherits the surrounding file as
context, leaving it intact. This way developers can safely leverage the
speed of AI generation without risking the stability of their existing
architecture. They maintain absolute control over the file's structure,
imports, and critical logic, effectively treating the AI as a
specialized contractor who is strictly forbidden from touching the
load-bearing walls of the application. This resolves the ``trust gap''
that currently prevents AI from being adopted in professional,
mission-critical codebases.

\textbf{Current limitations:} - No nesting support
(\texttt{\$\$\{\ \ \$\$\{...\}\$\$\ \ \}\$\$} fails) - Multi-line
handling requires careful parsing - No type constraints or formal
semantics

\textbf{2.1.1 Native Comments} Glupe introduces its own comment syntax
to annotate semantic logic without affecting the host language: -
\textbf{Line comments}: \texttt{\%\ This\ is\ a\ comment} -
\textbf{Block comments}:
\texttt{\%\{\ This\ is\ a\ multiline\ comment\ \}\%}

\subsubsection{2.2 Compilation Pipeline}\label{compilation-pipeline}

\begin{verbatim}
Input: Source file(s) with semantic blocks
Output: Executable binary or source code

Pipeline:
1. PREPROCESS: Resolve IMPORT directives (recursive module inclusion) and perform Directory Sucking
2. PARSE: Parse files using Flex/Bison frontend to extract semantic blocks and EXPORT directives, producing Glupe IR
3. GRAPH: Build Global Symbol Graph using Tree-sitter for targeted dependency injection
4. OPTIMIZE: If valid source without blocks → direct compile (bypass LLM)
5. GENERATE: Query LLM with context (surrounding code + errors + targeted injected dependencies)
6. VALIDATE: Compile with native toolchain and perform AST verification
7. ITERATE: On failure, retry with error feedback and semantic equivalence loop (max N attempts)
8. OUTPUT: Save binary or source file
\end{verbatim}

\textbf{Key insight}: Unlike traditional compilers that reject
incomplete programs, Glupe treats semantic blocks as \textbf{compilation
tasks} rather than syntax errors.

\subsubsection{2.3 Context Preservation}\label{context-preservation}

When generating code for a semantic block, Glupe provides the LLM with:

\begin{verbatim}
CONTEXT = {
    surrounding_code,      // Imports, type signatures, existing functions
    target_language,       // C++, Python, etc.
    previous_errors,       // Compilation failures from prior attempts
    user_instructions,     // Optional constraints via "*prompt"
    dependency_analysis    // Extracted #includes, imports
}
\end{verbatim}

Example:

\begin{Shaded}
\begin{Highlighting}[]
\PreprocessorTok{\#include }\ImportTok{\textless{}vector\textgreater{}}
\PreprocessorTok{\#include }\ImportTok{\textless{}algorithm\textgreater{}}

\DataTypeTok{int}\NormalTok{ main}\OperatorTok{()} \OperatorTok{\{}
    \BuiltInTok{std::}\NormalTok{vector}\OperatorTok{\textless{}}\DataTypeTok{int}\OperatorTok{\textgreater{}}\NormalTok{ data }\OperatorTok{=} \OperatorTok{\{}\DecValTok{5}\OperatorTok{,} \DecValTok{2}\OperatorTok{,} \DecValTok{8}\OperatorTok{,} \DecValTok{1}\OperatorTok{,} \DecValTok{9}\OperatorTok{\};}
    
    \ErrorTok{$$}\OperatorTok{\{}\NormalTok{sort the vector }\KeywordTok{and}\NormalTok{ remove duplicates}\OperatorTok{\}}\ErrorTok{$$}
    
    \ControlFlowTok{for}\OperatorTok{(}\KeywordTok{auto}\NormalTok{ x }\OperatorTok{:}\NormalTok{ data}\OperatorTok{)} \BuiltInTok{std::}\NormalTok{cout }\OperatorTok{\textless{}\textless{}}\NormalTok{ x }\OperatorTok{\textless{}\textless{}} \StringTok{" "}\OperatorTok{;}
\OperatorTok{\}}
\end{Highlighting}
\end{Shaded}

The LLM sees \texttt{\#include\ \textless{}algorithm\textgreater{}} and
knows to use \texttt{std::sort}, \texttt{std::unique} rather than
reimplementing sorting.

\subsubsection{2.4 Multi-File Projects: EXPORT
Directives}\label{multi-file-projects-export-directives}

For multi-file projects, Glupe introduces \texttt{EXPORT} blocks:

\begin{verbatim}
EXPORT: "filename.ext"
<file content with optional semantic blocks>
EXPORT: END
\end{verbatim}

Content inside \texttt{EXPORT} blocks is written to the specified file.
Content \textbf{outside} acts as project-level instructions without
being exported.

\textbf{Example:}

\begin{verbatim}
This project implements a calculator with separate interface/logic.
Use only standard library. No external dependencies.

EXPORT: "calculator.h"
#ifndef CALCULATOR_H
#define CALCULATOR_H
$${define Calculator class with add, subtract, multiply, divide methods}$$
#endif
EXPORT: END

EXPORT: "calculator.cpp"
#include "calculator.h"
$${implement Calculator methods}$$
EXPORT: END

EXPORT: "main.cpp"
#include <iostream>
#include "calculator.h"
$${create CLI interface for the calculator}$$
EXPORT: END
\end{verbatim}

The first two lines provide architectural context but don't appear in
any output file. This separates \textbf{project intent} from
\textbf{file structure}.

\textbf{Implementation:} - Simple line-by-line parser scanning for
\texttt{EXPORT:} markers - Filename extraction via regex for quoted
strings - Automatic directory creation for nested paths - Sequential or
parallel file generation.

\subsubsection{2.5 Advanced Semantics: Inheritance and Abstraction
(v5.8+)}\label{advanced-semantics-inheritance-and-abstraction-v5.8}

Glupe v5.8 introduced object-oriented concepts to semantic containers,
enabling architectural reuse and policy enforcement.

\textbf{1. Abstract Containers} Containers marked with \texttt{ABSTRACT}
do not generate code themselves but serve as templates for other
containers.

\begin{Shaded}
\begin{Highlighting}[]
\NormalTok{$$ABSTRACT secure\_api \{}
\NormalTok{    All database queries must use parameterized statements.}
\NormalTok{    All inputs must be sanitized.}
\NormalTok{    Use try{-}catch blocks for all network calls.}
\NormalTok{\}$$}
\end{Highlighting}
\end{Shaded}

\textbf{2. Container Inheritance} Containers can inherit intent from
parents using the \texttt{-\textgreater{}} operator. The child container
inherits the ``DNA'' (instructions) of the parent, automatically
applying constraints or logic.

\begin{Shaded}
\begin{Highlighting}[]
\NormalTok{$$ login\_handler {-}\textgreater{} secure\_api \{}
\NormalTok{    Implement the user login function.}
\NormalTok{\}$$}
\end{Highlighting}
\end{Shaded}

\textbf{3. Multiple Inheritance} Glupe supports mixing multiple
architectural traits:

\begin{Shaded}
\begin{Highlighting}[]
\NormalTok{$$ critical\_endpoint {-}\textgreater{} logging, metrics \{}
\NormalTok{    Handle the payment processing request.}
\NormalTok{\}$$}
\end{Highlighting}
\end{Shaded}

This solves the ``Drift'' problem: if the security policy changes in the
parent, every inheriting container updates automatically upon
recompilation.

\subsubsection{2.6 Prompt Algebra Basics}\label{prompt-algebra-basics}

When natural language intent is treated as a first-class data type, we
can impose strict mathematical topology on it. Instead of blindly
concatenating strings to send to an LLM, Glupe acts as an algebraic
``calculator,'' manipulating Semantic Nodes using structural operators
to compute a deterministic Final State of Intent.

Let I represent a Semantic Node (an Intent).

The core operations of Glupe's Intent Algebra include:

\textbf{1. Addition (Contextual Merging)} Combining two or more distinct
intents into a unified semantic space using Multiple Inheritance
(\texttt{-\textgreater{}\ A,\ B}). \emph{Formalization:}
\(\mathcal{I}_{child} = \mathcal{I}_{A} \oplus \mathcal{I}_{B}\)
\emph{Glupe Syntax:}
\texttt{\$\$\ C\ -\textgreater{}\ A,\ B\ \{\ build\ endpoint\ \}\$\$}
The compiler mathematically merges the constraints into a single,
unified prompt vector before transpilation. If the conditions are
contradictory, it throws a compile-time error.

\textbf{2. Subtraction (Constraint Overriding)} Removing specific
instruction members or negating inherited rules. In Glupe, explicit
contradiction in a child block or right-to-left precedence acts as a
subtractive override. \emph{Formalization:}
\(\mathcal{I}_{final} = \mathcal{I}_{parent} \ominus \mathcal{C}_{conflict}\)
\emph{Example:} A parent dictates ``Log all user data,'' while a child
overrides with ``Do not log passwords.'' The conflict is mathematically
bounded by the child's explicit subtraction.

\textbf{3. Multiplication (Functional Amplification)} Iterative
execution or mapping of one intent over another (Context Injection). By
passing a small, isolated intent \texttt{x} into a functional container
\texttt{F(x)}, the AI multiplies the effect of \texttt{x} across an
entire domain. \emph{Formalization:}
\(F(\mathcal{I}_x) = \mathcal{I}_x \otimes Iterator\) \emph{Glupe
Syntax:}
\texttt{\$\$\ process\_database(x)\ \{\ apply\ x\ to\ every\ column\ \}\$\$}

\textbf{4. The Identity Matrix (Abstract Variables)} A node that shapes
other nodes without possessing its own executable mass. It acts purely
as a mathematical limit (e.g.,
\texttt{\$ABSTRACT:\ MAX\_LATENCY\ -\textgreater{}\ 50ms}).
\emph{Formalization:}
\(\mathcal{I}_{final} = \mathcal{I}_{target} \cdot \mathbf{I}_{abstract}\)

Because Glupe resolves this algebra \emph{before} calling the LLM, it
evaluates the Abstract Syntax Tree (AST) of intent, substitutes
variables, merges bounds, and calculates the simplified deterministic
equation. If two intents demand mutually exclusive output states, their
sum mathematically resolves to an empty set (\(\bot\)), allowing the
compiler to throw a Compile-Time Error before any token is wasted on
inference.

\begin{center}\rule{0.5\linewidth}{0.5pt}\end{center}

\subsection{3. Reliability and Verification
Mechanisms}\label{reliability-and-verification-mechanisms}

LLMs are non-deterministic and prone to failures. Glupe implements
several mechanisms to improve reliability:

\subsubsection{3.1 Iterative Compilation with Error
Feedback}\label{iterative-compilation-with-error-feedback}

Rather than single-shot generation, Glupe uses a \textbf{retry loop with
accumulated error history}:

\begin{verbatim}
error_history = ""
for attempt in 1..MAX_RETRIES:
    code = LLM_generate(intent + context + error_history)
    result = native_compile(code)
    
    if result.success:
        return code
    
    error_history += "\n--- Attempt " + attempt + " failed ---\n"
    error_history += result.compiler_errors
    
    if is_fatal_error(result.errors):
        break
end
\end{verbatim}

Compiler error messages provide \textbf{ground truth feedback} that
guides the LLM toward valid solutions. We observe that most errors are
fixed within 2-4 iterations.

\subsubsection{3.2 Pattern-Based Failure
Detection}\label{pattern-based-failure-detection}

LLMs often attempt ``lazy'' solutions rather than native implementation:

\textbf{Common anti-patterns:} - Including
\texttt{\textless{}Python.h\textgreater{}} to embed Python in C++ -
Using \texttt{system("python\ script.py")} to shell out - Generating JNI
wrappers instead of Java code - Literal translation (writing Python
syntax in a C++ file)

\textbf{Detection strategy:}

\begin{Shaded}
\begin{Highlighting}[]
\DataTypeTok{bool}\NormalTok{ isFatalError}\OperatorTok{(}\AttributeTok{const}\NormalTok{ string}\OperatorTok{\&}\NormalTok{ compilerOutput}\OperatorTok{)} \OperatorTok{\{}
\NormalTok{    string lower }\OperatorTok{=}\NormalTok{ toLowerCase}\OperatorTok{(}\NormalTok{compilerOutput}\OperatorTok{);}
    
    \CommentTok{// Wrapper detection}
    \ControlFlowTok{if} \OperatorTok{(}\NormalTok{lower}\OperatorTok{.}\NormalTok{find}\OperatorTok{(}\StringTok{"python.h"}\OperatorTok{)} \OperatorTok{!=}\NormalTok{ string}\OperatorTok{::}\NormalTok{npos}\OperatorTok{)} \ControlFlowTok{return} \KeywordTok{true}\OperatorTok{;}
    \ControlFlowTok{if} \OperatorTok{(}\NormalTok{lower}\OperatorTok{.}\NormalTok{find}\OperatorTok{(}\StringTok{"jni.h"}\OperatorTok{)} \OperatorTok{!=}\NormalTok{ string}\OperatorTok{::}\NormalTok{npos}\OperatorTok{)} \ControlFlowTok{return} \KeywordTok{true}\OperatorTok{;}
    
    \CommentTok{// Literal translation detection}
    \ControlFlowTok{if} \OperatorTok{(}\NormalTok{lower}\OperatorTok{.}\NormalTok{find}\OperatorTok{(}\StringTok{"def "}\OperatorTok{)} \OperatorTok{!=}\NormalTok{ string}\OperatorTok{::}\NormalTok{npos}\OperatorTok{)} \ControlFlowTok{return} \KeywordTok{true}\OperatorTok{;}  \CommentTok{// Python function}
    \ControlFlowTok{if} \OperatorTok{(}\NormalTok{lower}\OperatorTok{.}\NormalTok{find}\OperatorTok{(}\StringTok{"print("}\OperatorTok{)} \OperatorTok{!=}\NormalTok{ string}\OperatorTok{::}\NormalTok{npos}\OperatorTok{)} \ControlFlowTok{return} \KeywordTok{true}\OperatorTok{;} \CommentTok{// Python print}
    
    \CommentTok{// Genuine missing dependencies}
    \ControlFlowTok{if} \OperatorTok{(}\NormalTok{lower}\OperatorTok{.}\NormalTok{find}\OperatorTok{(}\StringTok{"no such file"}\OperatorTok{)} \OperatorTok{!=}\NormalTok{ string}\OperatorTok{::}\NormalTok{npos}\OperatorTok{)} \ControlFlowTok{return} \KeywordTok{true}\OperatorTok{;}
    
    \ControlFlowTok{return} \KeywordTok{false}\OperatorTok{;}
\OperatorTok{\}}
\end{Highlighting}
\end{Shaded}

When detected, Glupe annotates the error message:

\begin{verbatim}
"FATAL: You attempted to include python.h. STOP. 
Rewrite using native C++ standard library only."
\end{verbatim}

This \textbf{trains the LLM} through reinforcement to avoid lazy
patterns.

\subsubsection{3.3 Preflight Dependency
Checking}\label{preflight-dependency-checking}

Before invoking the LLM, Glupe extracts all dependencies
(\texttt{\#include}, \texttt{import}, etc.) and verifies they exist:

\begin{Shaded}
\begin{Highlighting}[]
\NormalTok{dependencies }\OperatorTok{=}\NormalTok{ extract\_dependencies}\OperatorTok{(}\NormalTok{code}\OperatorTok{)}
\NormalTok{compile\_test\_file\_with\_dependencies}\OperatorTok{(}\NormalTok{dependencies}\OperatorTok{)}

\ControlFlowTok{if}\NormalTok{ compilation\_fails}\OperatorTok{:}
\NormalTok{    report\_missing\_dependencies}\OperatorTok{()}
\NormalTok{    abort\_before\_expensive\_LLM\_call}\OperatorTok{()}
\end{Highlighting}
\end{Shaded}

This prevents wasted API calls and tokens on doomed-to-fail generations.
\textgreater Note: as of version 5.8 this is only implemented for python
and c/c++

\subsubsection{3.4 Direct Compilation
Optimization}\label{direct-compilation-optimization}

\textbf{Key insight}: If input files are already valid source code in
the target language and contain no semantic blocks, bypass LLM entirely.

\begin{Shaded}
\begin{Highlighting}[]
\ControlFlowTok{if} \OperatorTok{(}\NormalTok{all\_files\_match\_target\_extension }\OperatorTok{\&\&}\NormalTok{ no\_semantic\_blocks}\OperatorTok{)} \OperatorTok{\{}
\NormalTok{    compile\_directly\_with\_native\_toolchain}\OperatorTok{();}
    \CommentTok{// Zero overhead for traditional workflows}
\OperatorTok{\}}
\end{Highlighting}
\end{Shaded}

This makes Glupe a \textbf{drop-in replacement} for native compilers
when no AI features are needed.

\subsubsection{3.5 AST-Driven Structural
Verification}\label{ast-driven-structural-verification}

LLMs are notorious for being ``lazy,'' often oversimplifying code when
asked to refactor or comment, replacing critical logic with
\texttt{//\ ...\ implementation\ goes\ here\ ...}. To counter this,
Glupe integrates Tree-sitter for Abstract Syntax Tree (AST)
manipulation.

During modification operations (like \texttt{fix} or \texttt{explain}),
Glupe mathematically compares the node count of structural milestones
(\texttt{function\_definition}, \texttt{if\_statement},
\texttt{for\_statement}, etc.) between the original code and the
AI-generated output. If the complexity ratio drops below a strict
threshold (e.g., \texttt{\textless{}\ 0.7}), indicating the AI dropped
significant structural logic, Glupe physically rejects the output and
forces the AI to try again with an explicit system alert. This
structural feedback loop forces non-deterministic AI models to maintain
algorithmic fidelity.

\subsubsection{3.6 Semantic Subtraction (The ``Omission Blindspot''
Guardrail)}\label{semantic-subtraction-the-omission-blindspot-guardrail}

A critical flaw in autonomous coding agents is their inability to detect
omissions. LLMs are decent at catching contradictions but terrible at
realizing what they forgot to implement from a specification. Standard
unit tests only verify \emph{if} written code runs, not \emph{how} it
was architected or if constraints were silently ignored.

Glupe introduces a mathematical verification gate called
\textbf{Semantic Subtraction (\(Spec — Impl\))}. 1. The developer
provides the Expected Intent (\texttt{spec.glp}). 2. Glupe
reverse-engineers the AI-generated code (\texttt{agent\_output.cpp})
back into a semantic blueprint (\texttt{actual\_impl.glp}) using
AST-driven refinement. 3. Glupe performs a two-way Set Subtraction
between both blueprints.

The \texttt{glupe\ audit} command mathematically evaluates the absence
of requested logic, catching missed requirements (Missing) and
architectural liberties (Hallucinated). Equipped with a
\texttt{-\/-ignore-scaffold} flag to forgive standard language
boilerplate, this provides a deterministic, CI/CD-ready guardrail that
can automatically block AI Pull Requests if they drift from the master
architectural intent.

\begin{center}\rule{0.5\linewidth}{0.5pt}\end{center}

\subsection{4. LLM Backend Architecture}\label{llm-backend-architecture}

\subsubsection{4.1 Multi-Provider Support}\label{multi-provider-support}

Glupe abstracts LLM providers through a unified interface:

\begin{Shaded}
\begin{Highlighting}[]
\NormalTok{string callAI}\OperatorTok{(}\NormalTok{string prompt}\OperatorTok{)} \OperatorTok{\{}
\NormalTok{    json body}\OperatorTok{;}
    
    \ControlFlowTok{if} \OperatorTok{(}\NormalTok{PROTOCOL }\OperatorTok{==} \StringTok{"google"}\OperatorTok{)} \OperatorTok{\{}
\NormalTok{        body}\OperatorTok{[}\StringTok{"contents"}\OperatorTok{][}\DecValTok{0}\OperatorTok{][}\StringTok{"parts"}\OperatorTok{][}\DecValTok{0}\OperatorTok{][}\StringTok{"text"}\OperatorTok{]} \OperatorTok{=}\NormalTok{ prompt}\OperatorTok{;}
\NormalTok{        url }\OperatorTok{+=} \StringTok{"?key="} \OperatorTok{+}\NormalTok{ API\_KEY}\OperatorTok{;}
    \OperatorTok{\}} 
    \ControlFlowTok{else} \ControlFlowTok{if} \OperatorTok{(}\NormalTok{PROTOCOL }\OperatorTok{==} \StringTok{"openai"}\OperatorTok{)} \OperatorTok{\{}
\NormalTok{        body}\OperatorTok{[}\StringTok{"messages"}\OperatorTok{][}\DecValTok{0}\OperatorTok{][}\StringTok{"role"}\OperatorTok{]} \OperatorTok{=} \StringTok{"user"}\OperatorTok{;}
\NormalTok{        body}\OperatorTok{[}\StringTok{"messages"}\OperatorTok{][}\DecValTok{0}\OperatorTok{][}\StringTok{"content"}\OperatorTok{]} \OperatorTok{=}\NormalTok{ prompt}\OperatorTok{;}
\NormalTok{        headers }\OperatorTok{+=} \StringTok{"Authorization: Bearer "} \OperatorTok{+}\NormalTok{ API\_KEY}\OperatorTok{;}
    \OperatorTok{\}}
    \ControlFlowTok{else} \OperatorTok{\{} \CommentTok{// ollama (local)}
\NormalTok{        body}\OperatorTok{[}\StringTok{"model"}\OperatorTok{]} \OperatorTok{=}\NormalTok{ MODEL\_ID}\OperatorTok{;}
\NormalTok{        body}\OperatorTok{[}\StringTok{"prompt"}\OperatorTok{]} \OperatorTok{=}\NormalTok{ prompt}\OperatorTok{;}
\NormalTok{        body}\OperatorTok{[}\StringTok{"stream"}\OperatorTok{]} \OperatorTok{=} \KeywordTok{false}\OperatorTok{;}
    \OperatorTok{\}}
    
\NormalTok{    response }\OperatorTok{=}\NormalTok{ curl\_post}\OperatorTok{(}\NormalTok{url}\OperatorTok{,}\NormalTok{ body}\OperatorTok{);}
    \ControlFlowTok{return}\NormalTok{ extract\_code}\OperatorTok{(}\NormalTok{response}\OperatorTok{);}
\OperatorTok{\}}
\end{Highlighting}
\end{Shaded}

\textbf{Supported backends:} - \textbf{Ollama} (default): Local models,
zero cost, privacy-preserving - \textbf{OpenAI}: GPT-4, GPT-5 via API -
\textbf{Google}: Gemini Pro, Flash via API - \textbf{Compatible APIs}:
Groq, together.ai, any OpenAI-compatible endpoint

Configuration via \texttt{config.json}:

\begin{Shaded}
\begin{Highlighting}[]
\FunctionTok{\{}
  \DataTypeTok{"local"}\FunctionTok{:} \FunctionTok{\{}
    \DataTypeTok{"model\_id"}\FunctionTok{:} \StringTok{"qwen2.5{-}coder:3b"}\FunctionTok{,}
    \DataTypeTok{"api\_url"}\FunctionTok{:} \StringTok{"http://localhost:11434/api/generate"}
  \FunctionTok{\},}
  \DataTypeTok{"cloud"}\FunctionTok{:} \FunctionTok{\{}
    \DataTypeTok{"protocol"}\FunctionTok{:} \StringTok{"openai"}\FunctionTok{,}
    \DataTypeTok{"model\_id"}\FunctionTok{:} \StringTok{"gpt{-}4{-}turbo"}\FunctionTok{,}
    \DataTypeTok{"api\_key"}\FunctionTok{:} \StringTok{"sk{-}..."}\FunctionTok{,}
    \DataTypeTok{"api\_url"}\FunctionTok{:} \StringTok{"https://api.openai.com/v1/chat/completions"}
  \FunctionTok{\},}
  \DataTypeTok{"max\_retries"}\FunctionTok{:} \DecValTok{15}
\end{Highlighting}
\end{Shaded}

This will open the browser in apifreellm.com website, where users can
access to free API keys.

Users can switch providers with flags:

\begin{Shaded}
\begin{Highlighting}[]
\ExtensionTok{glupe}\NormalTok{ hello.glp }\AttributeTok{{-}local}   \CommentTok{\# Use Ollama}
\ExtensionTok{glupe}\NormalTok{ hello.glp }\AttributeTok{{-}cloud}   \CommentTok{\# Use configured cloud provider}
\end{Highlighting}
\end{Shaded}

\subsubsection{4.2 Rate Limiting and
Backoff}\label{rate-limiting-and-backoff}

Cloud providers enforce rate limits. Glupe implements exponential
backoff:

\begin{Shaded}
\begin{Highlighting}[]
\ControlFlowTok{for} \OperatorTok{(}\DataTypeTok{int}\NormalTok{ attempt }\OperatorTok{=} \DecValTok{0}\OperatorTok{;}\NormalTok{ attempt }\OperatorTok{\textless{}} \DecValTok{3}\OperatorTok{;}\NormalTok{ attempt}\OperatorTok{++)} \OperatorTok{\{}
\NormalTok{    response }\OperatorTok{=}\NormalTok{ api\_call}\OperatorTok{(}\NormalTok{prompt}\OperatorTok{);}
    
    \ControlFlowTok{if} \OperatorTok{(}\NormalTok{response}\OperatorTok{.}\NormalTok{contains}\OperatorTok{(}\StringTok{"429"}\OperatorTok{)} \OperatorTok{||}\NormalTok{ response}\OperatorTok{.}\NormalTok{contains}\OperatorTok{(}\StringTok{"Rate limit"}\OperatorTok{))} \OperatorTok{\{}
\NormalTok{        wait\_seconds }\OperatorTok{=} \DecValTok{5} \OperatorTok{*} \OperatorTok{(}\NormalTok{attempt }\OperatorTok{+} \DecValTok{1}\OperatorTok{);}
\NormalTok{        sleep}\OperatorTok{(}\NormalTok{wait\_seconds}\OperatorTok{);}
        \ControlFlowTok{continue}\OperatorTok{;}
    \OperatorTok{\}}
    \ControlFlowTok{break}\OperatorTok{;}
\OperatorTok{\}}
\end{Highlighting}
\end{Shaded}

\subsubsection{4.3 Response Parsing}\label{response-parsing}

LLM responses vary by provider. Glupe extracts code blocks:

\begin{Shaded}
\begin{Highlighting}[]
\NormalTok{string extractCode}\OperatorTok{(}\NormalTok{string jsonResponse}\OperatorTok{)} \OperatorTok{\{}
\NormalTok{    json j }\OperatorTok{=}\NormalTok{ parse}\OperatorTok{(}\NormalTok{jsonResponse}\OperatorTok{);}
    
    \CommentTok{// OpenAI format}
    \ControlFlowTok{if} \OperatorTok{(}\NormalTok{j}\OperatorTok{.}\NormalTok{contains}\OperatorTok{(}\StringTok{"choices"}\OperatorTok{))}
\NormalTok{        text }\OperatorTok{=}\NormalTok{ j}\OperatorTok{[}\StringTok{"choices"}\OperatorTok{][}\DecValTok{0}\OperatorTok{][}\StringTok{"message"}\OperatorTok{][}\StringTok{"content"}\OperatorTok{];}
    
    \CommentTok{// Google format}
    \ControlFlowTok{else} \ControlFlowTok{if} \OperatorTok{(}\NormalTok{j}\OperatorTok{.}\NormalTok{contains}\OperatorTok{(}\StringTok{"candidates"}\OperatorTok{))}
\NormalTok{        text }\OperatorTok{=}\NormalTok{ j}\OperatorTok{[}\StringTok{"candidates"}\OperatorTok{][}\DecValTok{0}\OperatorTok{][}\StringTok{"content"}\OperatorTok{][}\StringTok{"parts"}\OperatorTok{][}\DecValTok{0}\OperatorTok{][}\StringTok{"text"}\OperatorTok{];}
    
    \CommentTok{// Ollama format}
    \ControlFlowTok{else} \ControlFlowTok{if} \OperatorTok{(}\NormalTok{j}\OperatorTok{.}\NormalTok{contains}\OperatorTok{(}\StringTok{"response"}\OperatorTok{))}
\NormalTok{        text }\OperatorTok{=}\NormalTok{ j}\OperatorTok{[}\StringTok{"response"}\OperatorTok{];}
    
    \CommentTok{// Strip markdown code fences}
    \ControlFlowTok{if} \OperatorTok{(}\NormalTok{text}\OperatorTok{.}\NormalTok{contains}\OperatorTok{(}\StringTok{"\textasciigrave{}\textasciigrave{}\textasciigrave{}"}\OperatorTok{))} \OperatorTok{\{}
        \ControlFlowTok{return}\NormalTok{ extract\_between\_fences}\OperatorTok{(}\NormalTok{text}\OperatorTok{);}
    \OperatorTok{\}}
    
    \ControlFlowTok{return}\NormalTok{ text}\OperatorTok{;}
\OperatorTok{\}}
\end{Highlighting}
\end{Shaded}

\begin{center}\rule{0.5\linewidth}{0.5pt}\end{center}

\subsection{5. Core Semantic Workflows}\label{core-semantic-workflows}

The Glupe paradigm enables several novel development workflows that are
not easily achievable with traditional tools.

\subsubsection{5.1 Semantic Refinement (Intent
Recovery)}\label{semantic-refinement-intent-recovery}

\begin{Shaded}
\begin{Highlighting}[]
\ExtensionTok{glupe}\NormalTok{ source.cpp }\AttributeTok{{-}refine}\NormalTok{ [{-}local }\KeywordTok{|} \ExtensionTok{{-}cloud]}
\end{Highlighting}
\end{Shaded}

\textbf{Workflow:} 1. Reads the source code file. 2. Prompts the LLM to
``semantically compress'' the implementation details into high-level
intent blocks (\texttt{\$\$\{...\}\$\$}), while preserving the
architectural skeleton (imports, class definitions, function
signatures). 3. Outputs a \texttt{.glp} file (e.g.,
\texttt{source.cpp.glp}).

\textbf{Implications:} - \textbf{Legacy Modernization}: Converts
``dead'' legacy code back into ``living'' intent. A developer can refine
a legacy C file into a \texttt{.glp} file, then recompile it targeting
Rust or Python. - \textbf{Intent Recovery}: Recovers the ``why'' behind
the code that is often lost in implementation details. -
\textbf{Codebase Compression}: Reduces the cognitive load for developers
reading the file, as they see the high-level logic first. - \textbf{It
allows for `no rot' software}: Unlike traditional code, intention does
not age, a program coded in \texttt{.glp} can be recompiled into the
same program in 10 or 20 years since it preserves the intent.

\subsubsection{5.2 Architectural Auditing (Semantic
Subtraction)}\label{architectural-auditing-semantic-subtraction}

\begin{Shaded}
\begin{Highlighting}[]
\ExtensionTok{glupe}\NormalTok{ audit spec.glp implementation\_refined.glp }\AttributeTok{{-}{-}ignore{-}scaffold}
\end{Highlighting}
\end{Shaded}

\textbf{Workflow:} 1. Read the expected architectural blueprint
(\texttt{spec.glp}). 2. Read the actual refined blueprint of the
generated code (\texttt{implementation\_refined.glp}). 3. Align
containers semantically (resolving AI naming drift). 4. Perform
mathematical Set Subtraction to generate a report of missing and
hallucinated features. 5. Return a standard exit code (\texttt{0} for
success, \texttt{1} for divergence) to act as a CI/CD pipeline
guardrail.

\textbf{Use cases:} - Mathematical verification of AI-generated Pull
Requests. - Test-Driven Architecture (TDA). - Ensuring strict adherence
to global security policies and architectures.

\textbf{Implications:} - \textbf{Zero-Trust AI Integration}:
Organizations no longer need to blindly trust that an LLM understood the
system architecture. They can mathematically prove that no constraints
were ignored and no unauthorized logic was hallucinated. -
\textbf{Automated Code Review}: Shifts the cognitive burden of reviewing
AI-generated code from human reviewers to a deterministic compiler gate.
- \textbf{Living Architecture Enforced}: The design blueprint
(\texttt{spec.glp}) ceases to be a stale wiki page and becomes a strict,
enforceable contract governing the repository.

\subsubsection{5.3 Targeted Implementation (Progressive
Enhancement)}\label{targeted-implementation-progressive-enhancement}

\begin{Shaded}
\begin{Highlighting}[]
\ExtensionTok{glupe}\NormalTok{ source.cpp }\AttributeTok{{-}fill}\NormalTok{ [{-}local }\KeywordTok{|} \ExtensionTok{{-}cloud]}
\end{Highlighting}
\end{Shaded}

\textbf{Workflow:} 1. Reads the source code file containing semantic
containers (\texttt{\$\$\{...\}\$\$}). 2. Prompts the LLM to generate
the implementation details exclusively for those containers, using the
surrounding code as context. 3. In-place replaces the containers with
the generated native code, leaving the rest of the file untouched.

\textbf{Use cases:} - \textbf{Boilerplate Elimination}: Write the
signatures and let the AI fill in the repetitive implementations. -
\textbf{Progressive Enhancement}: Iteratively fill in complex logic
blocks one by one while manually testing the surrounding architecture.

This calls for a spec-first driven coding paradigm.

\begin{center}\rule{0.5\linewidth}{0.5pt}\end{center}

\subsection{6. Language Support}\label{language-support}

\subsubsection{6.1 Language Profile
System}\label{language-profile-system}

Glupe supports 30+ languages through a declarative profile database:

\begin{Shaded}
\begin{Highlighting}[]
\KeywordTok{struct}\NormalTok{ LangProfile }\OperatorTok{\{}
\NormalTok{    string extension}\OperatorTok{;}        \CommentTok{// e.g., ".cpp"}
\NormalTok{    string versionCmd}\OperatorTok{;}       \CommentTok{// e.g., "g++ {-}{-}version"}
\NormalTok{    string buildCmd}\OperatorTok{;}         \CommentTok{// e.g., "g++ {-}std=c++17"}
    \DataTypeTok{bool}\NormalTok{ producesBinary}\OperatorTok{;}     \CommentTok{// true for compiled languages}
\OperatorTok{\};}

\NormalTok{map}\OperatorTok{\textless{}}\NormalTok{string}\OperatorTok{,}\NormalTok{ LangProfile}\OperatorTok{\textgreater{}}\NormalTok{ LANG\_DB }\OperatorTok{=} \OperatorTok{\{}
    \OperatorTok{\{}\StringTok{"cpp"}\OperatorTok{,} \OperatorTok{\{}\StringTok{".cpp"}\OperatorTok{,} \StringTok{"g++ {-}{-}version"}\OperatorTok{,} \StringTok{"g++ {-}std=gnu++17"}\OperatorTok{,} \KeywordTok{true}\OperatorTok{\}\},}
    \OperatorTok{\{}\StringTok{"py"}\OperatorTok{,}  \OperatorTok{\{}\StringTok{".py"}\OperatorTok{,}  \StringTok{"python {-}{-}version"}\OperatorTok{,} \StringTok{"python {-}m py\_compile"}\OperatorTok{,} \KeywordTok{false}\OperatorTok{\}\},}
    \OperatorTok{\{}\StringTok{"rust"}\OperatorTok{,\{}\StringTok{".rs"}\OperatorTok{,}  \StringTok{"rustc {-}{-}version"}\OperatorTok{,} \StringTok{"rustc"}\OperatorTok{,} \KeywordTok{true}\OperatorTok{\}\},}
    \CommentTok{// ... 30+ more}
\OperatorTok{\};}
\end{Highlighting}
\end{Shaded}

\textbf{Adding new language support:} 1. Add entry to
\texttt{config.json} with toolchain commands 2. No language-specific
parsing required 3. \textasciitilde5 lines of configuration

\textbf{Currently supported:} - \textbf{Compiled}: C, C++, Rust, Go,
Swift, Zig, Nim, Haskell, Dart - \textbf{Interpreted}: Python,
JavaScript, Ruby, PHP, Perl, Lua, Julia, R - \textbf{JVM}: Java, Kotlin,
Scala, Clojure - \textbf{Scripting}: Bash, PowerShell, Batch -
\textbf{Specialized}: TypeScript, LaTeX, SQL, HTML, CSS, Markdown

\subsubsection{6.2 Cross-Language
Transpilation}\label{cross-language-transpilation}

Same semantic description can target multiple languages:

\begin{Shaded}
\begin{Highlighting}[]
\ExtensionTok{glupe}\NormalTok{ program.Glupe }\AttributeTok{{-}cpp} \AttributeTok{{-}o}\NormalTok{ program.exe}
\ExtensionTok{glupe}\NormalTok{ program.Glupe }\AttributeTok{{-}py}  \AttributeTok{{-}o}\NormalTok{ program.py}
\ExtensionTok{glupe}\NormalTok{ program.Glupe }\AttributeTok{{-}rust} \AttributeTok{{-}o}\NormalTok{ program}
\end{Highlighting}
\end{Shaded}

\textbf{Use cases:} - Language exploration (``How would this look in
Rust?'') - Prototyping in Python, production in C++ - Educational
materials showing same algorithm in multiple languages

\subsubsection{6.3 Context assimilation}\label{context-assimilation}

Glupe supports this

\begin{verbatim}
Glupe script.py file.c scriptB.js -o app.exe -cpp 
\end{verbatim}

Glupe acts as a universal translator:

\begin{enumerate}
\def\labelenumi{\arabic{enumi}.}
\tightlist
\item
  Read Phase: It reads script.py (Python), file.c (C), and scriptB.js
  (JavaScript).
\item
  Context Extraction: It ignores the syntax differences and extracts the
  Logic and Intent (including any \[{ ... }\] blocks).
\item
  Harmonization: It builds a unified ``Project Context.'' It sees that
  script.py defines a data structure and file.c defines a sorting
  algorithm.
\item
  Transpilation: It melts down all that logic and re-casts it into the
  target language: C++ (-cpp).
\end{enumerate}

Output: It produces a single, unified native binary app.exe.

\begin{center}\rule{0.5\linewidth}{0.5pt}\end{center}

\subsection{7. Evaluation: Case Study on a Large-Scale, Real-World
Artifact}\label{evaluation-case-study-on-a-large-scale-real-world-artifact}

To evaluate the scalability and robustness of Semantic Metaprogramming,
we conducted a case study applying Glupe's \texttt{-refine} pipeline to
Niels Lohmann's \texttt{JSON\ for\ Modern\ C++}
(\texttt{nlohmann/json.hpp}). This single-header library comprises
approximately 25,000 lines of highly optimized, modern C++ code. It
serves as an ultimate stress test for any semantic compiler due to three
primary challenges:

\subsubsection{7.1 Challenge 1: The Preprocessor
Maze}\label{challenge-1-the-preprocessor-maze}

\texttt{nlohmann/json} relies heavily on preprocessor macros (e.g.,
\texttt{JSON\_HAS\_CPP\_17}, \texttt{\#ifdef}) to support multiple C++
standards simultaneously. Naive chunking algorithms based on
brace-counting fail catastrophically when encountering \texttt{\#ifdef}
blocks containing unmatched braces, fracturing the context.

\textbf{Solution:} By employing \textbf{AST-Driven Slicing} via
Tree-sitter, Glupe parses the C++ source into a structurally sound
Abstract Syntax Tree before any LLM interaction. It extracts discrete
logical units (e.g., \texttt{basic\_json} class methods) cleanly,
separating the preprocessor boilerplate from the algorithmic logic. This
ensures the LLM receives syntactically perfect ``thoughts'' to refine.

\subsubsection{7.2 Challenge 2: Deep Template Type
Aliasing}\label{challenge-2-deep-template-type-aliasing}

The library makes intense use of type aliasing
(\texttt{using\ json\ =\ basic\_json\textless{}...\textgreater{};}). A
traditional sliding-window context will forget these definitions by the
time it reaches line 15,000 to refine a serialization function, leading
to hallucinated intent.

\textbf{Solution:} Glupe utilizes a \textbf{Global Symbol Graph} to
perform Pass 1 extraction of all function signatures, class definitions,
and global variables. During Pass 2 refinement, it performs
\textbf{Targeted Dependency Injection}, providing the LLM with the exact
graph of dependencies required by the specific node, offering perfect
context without exceeding token limits.

\subsubsection{7.3 Challenge 3: Token Limit Exhaustion and Semantic
Loss}\label{challenge-3-token-limit-exhaustion-and-semantic-loss}

Feeding 25,000 lines of dense C++ (\textasciitilde150,000 tokens) into
an LLM often results in severe ``Lost in the Middle'' syndrome, where
the model skips over edge-case handling. For instance, highly precise
UTF-8 surrogate pair decoding might be lazily summarized by the LLM as
\texttt{\$\$\{\ decode\ utf8\ \}\$\$.}

\textbf{Solution:} Glupe implements an \textbf{In-Memory Semantic
Equivalence Verification Loop}. It transpiles the refined intent back to
C++ in memory and performs a structural AST diff against the original
function. If the generated mathematical nodes do not match the original
(e.g., missing surrogate math), Glupe automatically triggers an
auto-heal prompt, forcing the LLM to refine at a stricter granularity
level.

\subsubsection{7.4 Outcome}\label{outcome}

Through this AST-driven, mathematically verified refinement process,
Glupe successfully reverse-engineers complex, template-heavy C++ into a
declarative \texttt{.glp} blueprint, demonstrating that Semantic
Metaprogramming can scale to industrial-grade codebases while
maintaining structural fidelity.

\subsection{8. Related Work}\label{related-work}

\subsubsection{8.1 LLM-Assisted Programming
Tools}\label{llm-assisted-programming-tools}

\textbf{GitHub Copilot} \hyperref[references]{{[}1{]}}: IDE autocomplete
based on GPT Codex. Provides line/block suggestions but requires
accepting/rejecting each suggestion. Not integrated into compilation
pipeline.

\textbf{ChatGPT / Claude}: Full program generation through chat.
Excellent for greenfield development but disconnected from existing
codebases. Requires copy-paste workflow.

\textbf{Cursor}: AI-powered IDE with chat interface and codebase
awareness. Still operates through suggestions rather than compilable
syntax.

\textbf{Key difference}: Glupe treats semantic blocks as \textbf{part of
the source code} rather than external prompts. Generated code is
validated through native toolchains, not just ``accepted'' by the user.

\subsubsection{8.2 Intentional
Programming}\label{intentional-programming}

\textbf{Intentional Software (Simonyi, 2000s)}
\hyperref[references]{{[}2{]}}: Domain-specific notations compiled to
code through projectional editing. Required custom language workbenches.

\textbf{Key difference}: Glupe uses natural language interpreted by
LLMs, not custom DSLs. No special tooling required---semantic blocks
work in any text editor.

\subsubsection{8.3 Program Synthesis}\label{program-synthesis}

\textbf{Sketch} \hyperref[references]{{[}3{]}}: Constraint-based
synthesis from partial programs with holes.\\
\textbf{Rosette} \hyperref[references]{{[}4{]}}: Solver-aided synthesis
from formal specifications.

\textbf{Key difference}: These systems require formal specifications
(types, assertions, examples). Glupe accepts informal natural language.

\subsubsection{8.4 Natural Language
Programming}\label{natural-language-programming}

\textbf{AlphaCode} \hyperref[references]{{[}5{]}}: Generates competitive
programming solutions from problem descriptions.\\
\textbf{CodeGen} \hyperref[references]{{[}6{]}}: Multi-turn program
synthesis from conversational prompts.

\textbf{Key difference}: These systems generate complete programs from
scratch. Glupe enables \textbf{hybrid} programming where developers mix
concrete and abstract code.

\begin{center}\rule{0.5\linewidth}{0.5pt}\end{center}

\subsection{9. Limitations and Future
Work}\label{limitations-and-future-work}

\subsubsection{9.1 Current Limitations}\label{current-limitations}

\textbf{1. No formal semantics}\\
Semantic blocks lack type constraints. The LLM may generate code that
compiles but has incorrect semantics if the intent is ambiguous.

\emph{Example:} ``\$\({sort the array}\)'' could mean ascending,
descending, or by custom comparator.

\textbf{2. Non-determinism}\\
Same input may produce different outputs across runs due to LLM
sampling.

\emph{Mitigation:} Caching locks generated code for unchanged inputs.
Developers can ``freeze'' satisfactory outputs.

\textbf{3. No nesting support}\\
Semantic blocks cannot be nested:
\texttt{\$\$\{\ \ \$\$\{inner\}\$\$\ \ \}\$\$} fails to parse.

\textbf{4. Token limits}\\
Very large projects may exceed LLM context windows (typically 32K-128K
tokens).

\emph{Mitigation:} Series mode generates files sequentially with
accumulated context.

\textbf{6. Dependency on LLM quality}\\
Weak models (e.g., 1B parameter) produce poor code. System requires
competent models (3B+ parameters recommended).

\textbf{7. Security risks}\\
Generated code executes with user privileges. Malicious LLM output could
run arbitrary commands. Mitigation strategies like sandboxed execution
and static analysis are areas for future work.

\subsubsection{9.2 Future Directions}\label{future-directions}

\begin{enumerate}
\def\labelenumi{\arabic{enumi}.}
\tightlist
\item
  \textbf{Type-aware semantic blocks:}
\end{enumerate}

\begin{Shaded}
\begin{Highlighting}[]
\DataTypeTok{int} \ErrorTok{$$}\OperatorTok{\{}\NormalTok{compute factorial}\OperatorTok{\}}\ErrorTok{$$} \OperatorTok{(}\DataTypeTok{int}\NormalTok{ n}\OperatorTok{);}  \CommentTok{// LLM knows return type}
\end{Highlighting}
\end{Shaded}

Infer constraints from surrounding code (type signatures, variable
usage).

\textbf{2. Fine-tuning on project codebases:} Train domain-specific
models that understand project conventions, naming patterns,
architecture.

\textbf{3. Multi-modal input:}

\begin{Shaded}
\begin{Highlighting}[]
\ErrorTok{$$}\OperatorTok{\{}
    \OperatorTok{[}\NormalTok{diagram}\OperatorTok{.}\NormalTok{png}\OperatorTok{]}  \CommentTok{// Flowchart or architecture diagram}
\NormalTok{    implement }\KeywordTok{this}\NormalTok{ state machine}
\OperatorTok{\}}\ErrorTok{$$}
\end{Highlighting}
\end{Shaded}

\begin{center}\rule{0.5\linewidth}{0.5pt}\end{center}

\subsection{10. Conclusion}\label{conclusion}

We presented \textbf{Glupe}, a practical system for hybrid AI-assisted
programming where natural language intent (or any DSL) and concrete code
coexist as equal citizens. Our key insights:

\begin{enumerate}
\def\labelenumi{\arabic{enumi}.}
\tightlist
\item
  \textbf{Lightweight syntax} (\texttt{\$\$\{...\}\$\$}) integrates
  seamlessly with existing languages
\item
  \textbf{Iterative refinement} with compiler error feedback
  dramatically improves LLM reliability
\item
  \textbf{Pattern detection} catches common failure modes before wasted
  compilation attempts
\item
  \textbf{Multi-provider abstraction} enables local-first development
  with cloud fallback
\item
  \textbf{Developer tooling} (fix, explain, diff, sos) leverages same
  LLM backend for broader utility
\end{enumerate}

Glupe demonstrates that semantic programming is \textbf{practical today}
with current LLMs. As models improve, the balance may shift toward more
semantic, less syntactic programming---but the developer remains in
control, mixing precision where needed with abstraction where
beneficial.

\textbf{Semantic blocks are not a replacement for programming. They are
a new tool in the programmer's arsenal.}

\begin{center}\rule{0.5\linewidth}{0.5pt}\end{center}

\subsection{Acknowledgments}\label{acknowledgments}

Thanks to the open-source community for Ollama, nlohmann/json, and the
countless developers whose work enables local-first AI.

Thanks to Google for Gemini 3 VS Code extension, it was very useful all
the way.

\subsection{Thanks to Krzystof Dudek for useful feedback and
collaboration.}\label{thanks-to-krzystof-dudek-for-useful-feedback-and-collaboration.}

\subsection{References}\label{references}

{[}1{]} Chen, M., et al.~``Evaluating Large Language Models Trained on
Code.'' \emph{arXiv:2107.03374} (2021).

{[}2{]} Simonyi, C. ``The Death of Computer Languages, the Birth of
Intentional Programming.'' \emph{NATO Science Committee Conference}
(1995).

{[}3{]} Solar-Lezama, A., et al.~``Combinatorial Sketching for Finite
Programs.'' \emph{ASPLOS} (2006).

{[}4{]} Torlak, E., \& Bodik, R. ``A Lightweight Symbolic Virtual
Machine for Solver-Aided Host Languages.'' \emph{PLDI} (2014).

{[}5{]} Li, Y., et al.~``Competition-Level Code Generation with
AlphaCode.'' \emph{Science} 378.6624 (2022).

{[}6{]} Nijkamp, E., et al.~``CodeGen: An Open Large Language Model for
Code with Multi-Turn Program Synthesis.'' \emph{arXiv:2203.13474}
(2022).

{[}7{]} Roziere, B., et al.~``Code Llama: Open Foundation Models for
Code.'' \emph{arXiv:2308.12950} (2023).

{[}8{]} OpenAI. ``GPT-4 Technical Report.'' \emph{arXiv:2303.08774}
(2023).

\begin{center}\rule{0.5\linewidth}{0.5pt}\end{center}

\subsection{Appendix A: Command
Reference}\label{appendix-a-command-reference}

\subsubsection{Compilation}\label{compilation}

\begin{Shaded}
\begin{Highlighting}[]
\ExtensionTok{glupe}\NormalTok{ file.glp                      }\CommentTok{\# Compile with local LLM}
\ExtensionTok{glupe}\NormalTok{ file.glp }\AttributeTok{{-}cloud}               \CommentTok{\# Use cloud provider}
\ExtensionTok{glupe}\NormalTok{ file.glp }\AttributeTok{{-}cpp} \AttributeTok{{-}o}\NormalTok{ program      }\CommentTok{\# Target C++}
\ExtensionTok{glupe}\NormalTok{ file.glp }\AttributeTok{{-}py} \AttributeTok{{-}o}\NormalTok{ script.py     }\CommentTok{\# Target Python}
\ExtensionTok{glupe}\NormalTok{ file.glp }\AttributeTok{{-}make}                \CommentTok{\# Multi{-}file project mode}
\ExtensionTok{glupe}\NormalTok{ file.glp }\AttributeTok{{-}run}                 \CommentTok{\# Compile and execute}
\ExtensionTok{glupe}\NormalTok{ file.glp }\AttributeTok{{-}zig} \AttributeTok{{-}t}              \CommentTok{\# Transpile only (no binary)}
\end{Highlighting}
\end{Shaded}

\subsubsection{Developer Tools}\label{developer-tools}

\begin{Shaded}
\begin{Highlighting}[]
\ExtensionTok{glupe}\NormalTok{ fix file.cpp }\StringTok{"instruction"}   \CommentTok{\# AI{-}powered code repair}
\ExtensionTok{glupe}\NormalTok{ explain file.cpp Spanish     }\CommentTok{\# Generate documentation}
\ExtensionTok{glupe}\NormalTok{ diff v1.cpp v2.cpp           }\CommentTok{\# Semantic diff analysis}
\ExtensionTok{glupe}\NormalTok{ sos cpp }\StringTok{"error message"}      \CommentTok{\# Interactive help}
\ExtensionTok{glupe}\NormalTok{ file.cpp }\AttributeTok{{-}refine}             \CommentTok{\# Reverse engineer to .glp}
\ExtensionTok{glupe}\NormalTok{ file.cpp }\AttributeTok{{-}fill}               \CommentTok{\# Fill semantic containers in{-}place}
\end{Highlighting}
\end{Shaded}

\subsubsection{Configuration}\label{configuration}

\begin{Shaded}
\begin{Highlighting}[]
\ExtensionTok{glupe}\NormalTok{ config see                   }\CommentTok{\# View current config}
\ExtensionTok{glupe}\NormalTok{ config api{-}key YOUR\_KEY      }\CommentTok{\# Set cloud API key}
\ExtensionTok{glupe}\NormalTok{ config model{-}local           }\CommentTok{\# Select Ollama model (interactive)}
\ExtensionTok{glupe}\NormalTok{ config model{-}cloud gpt{-}4     }\CommentTok{\# Set cloud model}
\ExtensionTok{glupe}\NormalTok{ config max{-}retries 20        }\CommentTok{\# Set retry limit}
\end{Highlighting}
\end{Shaded}

\subsubsection{Utilities}\label{utilities}

\begin{Shaded}
\begin{Highlighting}[]
\ExtensionTok{glupe} \AttributeTok{{-}{-}init}                       \CommentTok{\# Create project template}
\ExtensionTok{glupe} \AttributeTok{{-}{-}version}                    \CommentTok{\# Show version}
\ExtensionTok{glupe} \AttributeTok{{-}{-}clean}                      \CommentTok{\# Remove temp files}
\ExtensionTok{glupe}\NormalTok{ clean cache                  }\CommentTok{\# Clear semantic cache}
\end{Highlighting}
\end{Shaded}

\begin{center}\rule{0.5\linewidth}{0.5pt}\end{center}

\subsection{Appendix B: Example Compilation
Trace}\label{appendix-b-example-compilation-trace}

\textbf{Input (\texttt{hello.glp}):}

\begin{verbatim}
EXPORT: "hello.cpp"
#include <iostream>

int main() {
    $${print hello world 5 times}$$
    return 0;
}
EXPORT: END
\end{verbatim}

\textbf{Compilation Log:}

\begin{verbatim}
[CHECK] Toolchain for C++...
   [OK] Ready.
[INFO] 'EXPORT:' directive detected. Auto-enabling Architect Mode (-make).
[CHECK] Verifying dependencies locally...
   [OK] Dependencies verified.
   [Pass 1] Architecting Project...
[EXPORT] Writing to hello.cpp...
   Verifying...
BUILD SUCCESSFUL: hello.cpp
   [Source]: hello.cpp
\end{verbatim}

\textbf{Output (\texttt{hello.cpp}):}

\begin{Shaded}
\begin{Highlighting}[]
\PreprocessorTok{\#include }\ImportTok{\textless{}iostream\textgreater{}}

\DataTypeTok{int}\NormalTok{ main}\OperatorTok{()} \OperatorTok{\{}
    \ControlFlowTok{for}\OperatorTok{(}\DataTypeTok{int}\NormalTok{ i }\OperatorTok{=} \DecValTok{0}\OperatorTok{;}\NormalTok{ i }\OperatorTok{\textless{}} \DecValTok{5}\OperatorTok{;}\NormalTok{ i}\OperatorTok{++)} \OperatorTok{\{}
        \BuiltInTok{std::}\NormalTok{cout }\OperatorTok{\textless{}\textless{}} \StringTok{"Hello World"} \OperatorTok{\textless{}\textless{}} \BuiltInTok{std::}\NormalTok{endl}\OperatorTok{;}
    \OperatorTok{\}}
    \ControlFlowTok{return} \DecValTok{0}\OperatorTok{;}
\OperatorTok{\}}
\end{Highlighting}
\end{Shaded}

\begin{center}\rule{0.5\linewidth}{0.5pt}\end{center}

\subsection{Appendix C: Architecture
Diagram}\label{appendix-c-architecture-diagram}

\begin{verbatim}
┌─────────────────────────────────────────────────────────┐
│                    Glupe COMPILER                        │
│                                                         │
│  Input: .Glupe files with semantic blocks                │
│         ↓                                               │
│  ┌──────────────────────────────────────────┐           │
│  │ 1. PREPROCESSOR                          │           │
│  │    - Resolve IMPORT: directives          │           │
│  │    - Detect EXPORT: blocks               │           │
│  │    - Check for semantic blocks           │           │
│  └────────────┬─────────────────────────────┘           │
│               ↓                                         │
│  ┌──────────────────────────────────────────┐           │
│  │ 2. OPTIMIZATION CHECK                    │           │
│  │    - If valid source + no blocks         │           │
│  │      → Direct compilation (bypass LLM)   │           │
│  └────────────┬─────────────────────────────┘           │
│               ↓                                         │
│  ┌──────────────────────────────────────────┐           │
│  │ 3. CONTEXT ASSEMBLY                      │           │
│  │    - Extract surrounding code            │           │
│  │    - Analyze dependencies                │           │
│  │    - Prepare LLM prompt                  │           │
│  └────────────┬─────────────────────────────┘           │
│               ↓                                         │
│  ┌──────────────────────────────────────────┐           │
│  │ 4. LLM BACKEND (Multi-Provider)          │           │
│  │    ┌─────────────────────────────────┐   │           │
│  │    │ Ollama (Local)                  │   │           │
│  │    │ OpenAI (Cloud)                  │   │           │
│  │    │ Google (Cloud)                  │   │           │
│  │    └─────────────────────────────────┘   │           │
│  └────────────┬─────────────────────────────┘           │
│               ↓                                         │
│  ┌──────────────────────────────────────────┐           │
│  │ 5. VALIDATION & ITERATION                │           │
│  │    - Compile with native toolchain       │           │
│  │    - Detect errors/anti-patterns         │           │
│  │    - Retry with error feedback           │           │
│  │    - Max 15 attempts                     │           │
│  └────────────┬─────────────────────────────┘           │
│               ↓                                         │
│  ┌──────────────────────────────────────────┐           │
│  │ 6. OUTPUT                                │           │
│  │    - Binary (if compiled language)       │           │
│  │    - Source (if interpreted/requested)   │           │
│  │    - Multi-file exports (if -make mode)  │           │
│  └──────────────────────────────────────────┘           │
│                                                         │
└─────────────────────────────────────────────────────────┘
\end{verbatim}

\begin{center}\rule{0.5\linewidth}{0.5pt}\end{center}

\textbf{Project Repository:}
\texttt{https://github.com/alonsovm44/glupe}\\
\textbf{License:} MIT\\
\textbf{Contact:} alonsovm443@outlook.com

\end{document}
